\documentclass{article}

% packages
\usepackage{amsmath, amsthm, thmtools, amsfonts, amssymb, luacode, catchfile, tikzducks, hyperref, ifthen}
\ifcsname c@kobocompile\endcsname
	\usepackage[a5paper, total={1072pt, 1448pt}, margin=10pt, includeheadfoot]{geometry} % set page margins
\else
	\usepackage[a4paper, margin=50pt, includeheadfoot]{geometry}
\fi
\usepackage[shortlabels]{enumitem}
\usepackage[skip=3pt, indent=0pt]{parskip}

% language
\usepackage[bidi=basic, layout=tabular, provide=*]{babel}
\ifcsname c@english\endcsname
	\babelprovide[main, import]{english}
\else
	\babelprovide[main, import]{hebrew}
	\babelprovide{rl}
\fi
%\babelfont{rm}{Libertinus Serif}
\babelfont{rm}[Renderer=Harfbuzz]{Libertinus Serif}
\babelfont{sf}{Libertinus Sans}
\babelfont{tt}{Libertinus Mono}

% style
\AddToHook{cmd/section/before}{\clearpage}	% Add line break before section
\linespread{1.3}
\setcounter{secnumdepth}{0}		% Remove default number tags from sections, this won't do well with theorems
\AtBeginDocument{\setlength{\belowdisplayskip}{3pt}}
\AtBeginDocument{\setlength{\abovedisplayskip}{3pt}}
\graphicspath{ {../images/} }

% operators
\DeclareMathOperator\cis{cis}
\DeclareMathOperator\Sp{Sp}
\DeclareMathOperator\tr{tr}
\DeclareMathOperator\im{Im}
\DeclareMathOperator\re{Re}
\DeclareMathOperator\diag{diag}
\DeclareMathOperator*\lowlim{\underline{lim}}
\DeclareMathOperator*\uplim{\overline{lim}}
\DeclareMathOperator\rng{rng}
\DeclareMathOperator\Sym{Sym}
\DeclareMathOperator\Arg{Arg}
\DeclareMathOperator\Log{Log}
\DeclareMathOperator\dom{dom}
\DeclareMathOperator\supp{Supp}
\DeclareMathOperator\var{Var}
\DeclareMathOperator\cov{Cov}

% commands
%\renewcommand\qedsymbol{\textbf{מש''ל}}
%\renewcommand\qedsymbol{\fbox{\emoji{lizard}}}
\newcommand{\Aa}[0]{\mathcal{A}}
\newcommand{\Bb}[0]{\mathcal{B}}
\newcommand{\CC}[0]{\mathbb{C}}
\newcommand{\Cc}[0]{\mathcal{C}}
\newcommand{\EE}[0]{\mathbb{E}}
\newcommand{\FF}[0]{\mathbb{F}}
\newcommand{\Ff}[0]{\mathcal{F}}
\newcommand{\Ii}[0]{\mathcal{I}}
\newcommand{\Gg}[0]{\mathcal{G}}
\newcommand{\Ll}[0]{\mathcal{L}}
\newcommand{\Mm}[0]{\mathcal{M}}
\newcommand{\NN}[0]{\mathbb{N}}
\newcommand{\Nn}[0]{\mathcal{N}}
\newcommand{\PP}[0]{\mathbb{P}}
\newcommand{\Pp}[0]{\mathcal{P}}
\newcommand{\QQ}[0]{\mathbb{Q}}
\newcommand{\RR}[0]{\mathbb{R}}
\newcommand{\Rr}[0]{\mathcal{R}}
\newcommand{\Ss}[0]{\mathcal{S}}
\newcommand{\TT}[0]{\mathbb{T}}
\newcommand{\Uu}[0]{\mathcal{U}}
\newcommand{\Vv}[0]{\mathcal{V}}
\newcommand{\Ww}[0]{\mathcal{W}}
\newcommand{\ZZ}[0]{\mathbb{Z}}
\newcommand{\acts}[0]{\circlearrowright}
\newcommand{\explain}[2] {
	\begin{flalign*}
		 && \text{#2} && \text{#1}
	\end{flalign*}
}
\newcommand{\maketitleprint}[0]{ \begin{center}
	%\begin{tikzpicture}[scale=3]
	%	\duck[graduate=gray!20!black, tassel=red!70!black]
	%\end{tikzpicture}	
	\includegraphics[width=6cm]{cover}
\end{center}
}

% theorem commands
\newtheoremstyle{c_remark}
	{}	% Space above
	{}	% Space below
	{}% Body font
	{}	% Indent amount
	{\bfseries}	% Theorem head font
	{}	% Punctuation after theorem head
	{.5em}	% Space after theorem head
	{\thmname{#1}\thmnumber{ #2}\thmnote{ \normalfont{\text{(#3)}}}}	% head content
\newtheoremstyle{c_definition}
	{3pt}	% Space above
	{3pt}	% Space below
	{}% Body font
	{}	% Indent amount
	{\bfseries}	% Theorem head font
	{}	% Punctuation after theorem head
	{.5em}	% Space after theorem head
	{\thmname{#1}\thmnumber{ #2}\thmnote{ \normalfont{\text{(#3)}}}}	% head content
\newtheoremstyle{c_plain}
	{3pt}	% Space above
	{3pt}	% Space below
	{\itshape}% Body font
	{}	% Indent amount
	{\bfseries}	% Theorem head font
	{}	% Punctuation after theorem head
	{.5em}	% Space after theorem head
	{\thmname{#1}\thmnumber{ #2}\thmnote{ \text{(#3)}}}	% head content

\ifcsname c@english\endcsname
	\theoremstyle{plain}
	\newtheorem{theorem}{Theorem}[section]
	\newtheorem{lemma}[theorem]{Lemma}
	\newtheorem{proposition}[theorem]{Proposition}
	\newtheorem*{proposition*}{Proposition}
	%\newtheorem{corollary}[theorem]{אין חלופה עברית}

	\theoremstyle{definition}
	\newtheorem{definition}[theorem]{Definition}
	\newtheorem*{definition*}{Definition}
	\newtheorem{example}{Example}[section]
	\newtheorem{exercise}{Exercise}[section]

	\theoremstyle{remark}
	\newtheorem*{remark}{Remark}
	\newtheorem*{solution}{Solution}
	\newtheorem{conclusion}[theorem]{Conclusion}
	\newtheorem{notation}[theorem]{Notation}
\else
	\theoremstyle{c_plain}
	\newtheorem{theorem}{משפט}[section]
	\newtheorem{lemma}[theorem]{למה}
	\newtheorem{proposition}[theorem]{טענה}
	\newtheorem*{proposition*}{טענה}
	%\newtheorem{corollary}[theorem]{אין חלופה עברית}

	\theoremstyle{c_definition}
	\newtheorem{definition}[theorem]{הגדרה}
	\newtheorem*{definition*}{הגדרה}
	\newtheorem{example}{דוגמה}[section]
	\newtheorem{exercise}{תרגיל}[section]

	\theoremstyle{c_remark}
	\newtheorem*{remark}{הערה}
	\newtheorem*{solution}{פתרון}
	\newtheorem{conclusion}[theorem]{מסקנה}
	\newtheorem{notation}[theorem]{סימון}
\fi

% Questions related commands
\newcounter{question}
\setcounter{question}{1}
\newcounter{sub_question}
\setcounter{sub_question}{1}

\ifcsname c@english\endcsname
	\newcommand{\question}[1][0]{
		\ifthenelse{#1 = 0}{}{\setcounter{question}{#1}}
		\section{Question \arabic{question}}
		\addtocounter{question}{1}
		\setcounter{sub_question}{1}
	}

	\newcommand{\subquestion}[1][0]{
		\ifthenelse{#1 = 0}{}{\setcounter{sub_question}{#1}}
		\subsection{Part \alph{sub_question}}
		\addtocounter{sub_question}{1}
	}
\else
	\newcommand{\question}[1][0]{
		\ifthenelse{#1 = 0}{}{\setcounter{question}{#1}}
		\section{שאלה \arabic{question}}
		\addtocounter{question}{1}
		\setcounter{sub_question}{1}
	}

	\newcommand{\subquestion}[1][0]{
		\ifthenelse{#1 = 0}{}{\setcounter{sub_question}{#1}}
		\subsection{סעיף \localecounter{letters.gershayim}{sub_question}}
		\addtocounter{sub_question}{1}
	}
\fi

% import lua and start of document
\directlua{common = require ('../common')}

\GetEnv{AUTHOR}

% headers
\author{\AUTHOR}
\date\today

\title{פתרון מטלה 01 --- אנליזה על יריעות, 80426}
\usepackage{pgfplots}
\pgfplotsset{compat=1.18}

\begin{document}
\maketitle
\maketitleprint{}

\question{}
תהי המסילה $\gamma : [0, 2] \to \RR^2$ המוגדרת על־ידי,
\[
	\gamma(t) = (t^2, t^3)
\]

\subquestion{}
נחשב את האורך של $\gamma$.
\begin{solution}
	נבחין כי $\dot{\gamma}(t) = (2t, 3t^2)$ מחישוב ישיר, וכן מהגדרה,
	\[
		l(\gamma)
		= \int_\gamma 1\ dl
		= \int_0^2 \lVert \dot{\gamma}(t) \rVert\ dt
		= \int_0^2 \sqrt{4t^2 + 9t^4}\ dt
		= \int_0^2 t \sqrt{4 + 9t^2}\ dt
	\]
	נשתמש בכלל ההצבה עבור $u = t^2, du = 2t dt$ ונקבל
	\[
		\int_0^2 t \sqrt{4 + 9t^2}\ dt
		= \frac{1}{2} \int_0^4 {(4 + 9u)}^\frac{1}{2}\ dt
		= \left. \frac{1}{2} {(4 + 9u)}^\frac{3}{2} \cdot \frac{1}{9} \cdot \frac{2}{3} \right\rvert_{u = 0}^{u = 2}
		= \frac{1}{3} ({40}^\frac{3}{2} - 4^\frac{3}{2})
	\]
\end{solution}

\subquestion{}
נחשב את הישר המשיק ואת משיק היחידה ל־$\gamma$ ב־$t = 1$.
\begin{solution}
	לפי הגדרת הישר המשיק נחשב ונקבל,
	\[
		l
		= \Sp\{ \dot{\gamma}(1) \}
		= \Sp\{ (2, 3) \}
		= \{ (2, 3)t \mid t \in \RR \}
	\]
	וכן משיק היחידה הוא
	\[
		\vec{T}_\gamma(1)
		= \frac{\dot{\gamma}(1)}{\lVert \dot{\gamma(1)} \rVert}
		= \frac{(2, 3)}{\sqrt{13}}
	\]
\end{solution}

\question{}
עבור המסילות הבאות נצייר ונחשב את האינטגרל המסילתי של $\gamma : [0, 1] \to \RR^n$ עם הפונקציות הנתונות.

\subquestion{}
נגדיר $\gamma(t) = (t^2, t^2 - t)$ ונחשב את $\int_\gamma F\ d\gamma$ עבור $F(x, y) = (x^2 y, y - 3x)$.
\begin{solution}
	עבור הציור נרצה לבצע רפרצטריזציה של $\gamma$ כך ש־$\gamma \circ \mu = (t, f(t))$, נבחין כי עבור $\mu(t) = \sqrt{t}$ נקבל $f(t) = t - \sqrt{t}$ ב־$[0, 1]$ ולכן נוכל לעבור לציור,
	\begin{otherlanguage}{english}
		\begin{center}
			\begin{tikzpicture}
				\begin{axis}[ymin=-0.5,ymax=0.5,xmin=0,xmax=1, samples=300]
					\addplot[blue, ultra thick] (x,{x - sqrt(x)});
				\end{axis}
			\end{tikzpicture}
		\end{center}
	\end{otherlanguage}
	ונעבור לחישוב האינטגרל,
	\begin{align*}
		\int_\gamma F\ d\gamma
		& = \int_0^1 F(\gamma(t)) \cdot \dot{\gamma}(t)\ dt \\
		& = \int_0^1 (t^4 (t^2 - t), t^2 - t - 3t^2) \cdot (2t, 2t - 1)\ dt \\
		& = \int_0^1 2t^5 (t^2 - t) + (2t - 1) (-t - 2t^2)\ dt \\
		& = \int_0^1 2t^7 - 2t^6 - 4t^3 - t\ dt \\
		& = \left. \frac{2}{8} t^8 - \frac{2}{7} t^7 - t^4 - \frac{1}{2} t^2 \right\rvert_{t = 0}^{t = 1} \\
		& = \frac{1}{4} - \frac{2}{7} - 1 - \frac{1}{2} - 0
	\end{align*}
\end{solution}

\subquestion{}
נגדיר $\gamma(t) = (t, \cos t, \sin t)$ וכן $F(x, y, z) = (0, -z, y)$.
\begin{solution}
	נצייר,
	\begin{otherlanguage}{english}
		\begin{center}
			\begin{tikzpicture}
				\begin{axis}
				\addplot3[
					surf,
					domain=0:1,
					] (x,{cos(x)},{sin(x)});
				\end{axis}
			\end{tikzpicture}
		\end{center}
	\end{otherlanguage}
	נעבור לחישוב האינטגרל,
	\[
		\int_\gamma F\ d\gamma
		= \int_0^1 F(\gamma(t)) \cdot \dot{\gamma}(t)\ dt
		= \int_0^1 (0, -\sin t, \cos t) \cdot (1, -\sin t, \cos t)\ dt
		= \int_0^1 \sin^2 t + \cos^2 t\ dt
		= \int_0^1 1\ dt
		= 1
	\]
\end{solution}

\subquestion{}
נגדיר $\gamma(t) = (2t - 1, \sqrt{1 - {(2t - 1)}^2})$, ואת הפונקציה $f(x, y) = x y^4$.
\begin{solution}
	נתחיל בציור,
	\begin{otherlanguage}{english}
		\begin{center}
			\begin{tikzpicture}
				\begin{axis}[
					%ymin=-0.5,ymax=0.5,
					%xmin=0,xmax=1,
					samples=300
					]
					\addplot[black, thick] ({2 * x - 1},{sqrt(1 - (2 * x - 1)^2)});
				\end{axis}
			\end{tikzpicture}
		\end{center}
	\end{otherlanguage}
	ונעבור לחישוב האינטגרל.
	נבצע רפרמטריזציה יחד עם $\mu : [-1, 1] \to [0, 1]$ על־ידי $\mu(t) = \frac{t + 1}{2}$ ונקבל $(\gamma \circ \mu)(t) = (t, \sqrt{1 - t^2})$,
	ולאחר מכן נקבל $\bar{\gamma} : [0, \pi] \to \RR^2$ המוגדרת על־ידי $\bar{\gamma}(t) = (\cos t, \sin t)$.
	בנוסף גם $\dot{\bar{\gamma}}(t) = (-\sin t, \cos t)$ וכן $f(\bar{\gamma}(t)) = \cos(t) \sin^4(t)$, לכן
	\[
		\int_{\bar{\gamma}} f\ ds
		= \int_0^\pi \cos(t) \sin^4(t) \sqrt{\sin^2(t) + \cos^2(t)}\ dt
		= \left. \frac{1}{5} \sin^5(t) \right\rvert_{t = 0}^{t = \pi}
		= 0
	\]
\end{solution}

\question{}
תהי מסילה רגולרית $\gamma : [a, b] \to U$ ותהי $\bar{\gamma}$ המסילה נורמלית שלה. \\
נוכיח שכמות האורך המכוסה על־ידי $\bar{\gamma}$ בכל קטע היא האורך של הקטע, כלומר לכל $s_1 < s_2 \in [0, l(\gamma)]$ מתקיים $l(\bar{\gamma} \mid_{[s_1, s_2]}) = s_2 - s_1$.
\begin{proof}
	
\end{proof}

\end{document}
